% Lösungen Einheit 3
\einheitenkopf[le3][3 Sachaufgaben]{Einheit 3 von 4 · Sachaufgaben}
\erg{40}{a) 9 \quad b) beide \quad c) 6 \quad d) 12 \quad e) 7}
\erg{41}{a) 13 und −13 \quad b) 20 und −20 \quad c) x² = 64; 8 und −8 \quad d) x² = 100; 10 und −10}
\erg{42}{a) x² + 15 = 96, x² = 81; 9 und −9 \quad b) 3x² = 75, x² = 25; 5 und −5 \quad c) x · (x + 1) = 56, x² + x − 56 = 0; $x_1$ = 7, $x_2$ = −8; Zahlen 7 und 8 oder −8 und −7 \quad d) x · (x + 1) = 132, x² + x − 132 = 0; $x_1$ = 11, $x_2$ = −12; nur 11 und 12, weil natürliche Zahlen nicht negativ sind}
\erg{43}{a) x · (x + 3) = 40 \quad b) x² + 3x − 40 = 0; −1,5 $\pm$ $\surd$(2,25 + 40) = −1,5 $\pm$ 6,5; $x_1$ = 5, $x_2$ = −8 \quad c) nur x = 5, eine Länge ist nicht negativ \quad d) Breite 5 m, Länge 8 m; Probe: 5 m · 8 m = 40 m²}
\erg{44}{a) (x + 4)² = 196; x + 4 = 14 oder x + 4 = −14; $x_1$ = 10, $x_2$ = −18 (entfällt); Seite des Fotos 10 cm \quad b) (x + 4) · (x − 2) = 55; x² + 2x − 63 = 0; −1 $\pm$ 8; $x_1$ = 7, $x_2$ = −9 (entfällt); Seite des Gartens 7 m}
\erg{45}{a) Tim nimmt die negative Lösung als Länge: nur x = 3; Breite 3 cm, Länge 8 cm \quad b) Mia rechnet mit dem Umfang statt mit der Fläche: x · (x + 2) = 48, x² + 2x − 48 = 0; $x_1$ = 6, $x_2$ = −8 (entfällt); Breite 6 cm, Länge 8 cm}
\erg{46}{a) Eine Länge kann nicht negativ sein. \quad b) 15 und −15 ergeben quadriert beide 225, und das Rätsel verlangt keine positive Zahl.}
